\section{Principle VIII: Heredity}

\begin{quotex}
Where there is no vision, the people perish: but he that keeps the law, happy is he. \flright{\textsc{Proverbs 29:18}}

In this insightful section, Maurras shows the importance of heredity, both in regard to progeny, but also with respect to the transmission of the ``fruits" of a society through its bloodline. He points out that it is the society of those with faith who are most likely to understand this, and that it is due to God's very law: Be fruitful and multiply. Those angst-filled existentialists can fulfill this task only paradoxically by making it a sacred task. As faith declines, men see no need to preserve and pass on their heritage. 

\end{quotex}
Nature proceeds most commonly by imitation and repetition: but it also has innovations, caprices, initiatives. In that case it acts with an extreme vigor.

When it lets a soldier be born in the house of peaceable magistrates, or a marine in the descent of wine growers, the new vocation is marked rather strongly, it is served by a rather firm will so that all resistances are broken. But these resistances, these difficulties have some good. These hardships are natural tests, letting the predestined strong pass, but rejecting the others in the hereditary condition most suitable to them, because nature assures them defense and protection.

As the means of action towards a future, heredity is the most direct and the simplest of all.

Its general utility results from the destiny of the generator gifted with reason, who reproduces himself before dying.

Human life would be shamefully short if nature had not furnished societies with a procedure that transmits the fruits of its works through the blood.

The impersonal passion, whose deposit man holds, acts only through him in order to get to others, but delegates to the children which he procreates a power over the goods that he likewise procreated, very often for them. The inert posterity which comes from his hand will be vivified by his living posterity. When his great sons begin to enrich his heritage, all workers a little amorous of his work feel with some truth that he is going to conquer death twice. The power of bequeathing his remains gives to the activity of a well-filled life the highest natural laurels.

Let us note that the Christian societies of the Middle Age, penetrated by the supernatural sentiment of a future life, always proved themselves extremely sensitive to the terrestrial reward of the father in his sons. They sang, with all their soul, ``Abraham et semini ejus in secula" [``to Abraham and to his sons forever"].

Our different royal or imperial races arise out of nations fervently convinced of the reality of the kingdom of heaven: how would less believing peoples pay less attention to the carnal wish of hereditary duration? It is their only defense against time; they have only that anchor to throw on the abyss of the future. In a pinch, hereditary ambitions should have weakened through a sudden flight of celestial hopes; would the opposite be understood?

We can reason this way:

If some God hidden in the secret of hearts or soaring on the interplanetary abyss assists, immobile and mute, ardent and all powerful, the development of the efforts of humanity, it is His own law that He confirms in things and men; He can only bless its increased effects.

But if the spaces are empty and if the human heart is itself not assisted by any ``internal consolation", all the happiness of the being and all the benefits of life appear more exposed to the erosion of time and the blow of death, their tradition, their transmission then seems more precious in its solitary immensity to perpetual destitution. All means of saving or prolonging the shaky personal effort becomes more sacred perhaps! The thought threatened is attached more closely to the philosophy of order and the knowledge of the laws of its preservation. Should this order succumb, the believers keep refuge in the divine City: he who no longer believes undergoes the catastrophe of everything that his dream vied with death.



\flrightit{Posted on 2012-12-09 by Cologero }
